\chapter{\MakeUppercase{System Design}}
\justifying
\section{System Architecture}
\paragraph{} The adaptive gameplay system adopts a four-layer architecture, with each layer responsible for a specific part of the overall system. This modular approach ensures a clean separation of concerns and allows each component to be developed and tested independently.

\subsection{Gameplay Layer (Pygame)}
\paragraph{} The Gameplay Layer is implemented in \texttt{game\_3d.py} and handles all game mechanics:
\begin{itemize}
	\item \textbf{Player Management:} Handles spaceship movement (WASD/Arrow keys), shooting (Spacebar), health, and shield management.
	\item \textbf{Enemy System:} Controls enemy ship spawning, movement, health, and attack behavior. Parameters such as enemy speed, health, spawn delay, and attack intensity are presented as the adjustable variables being controlled by the \ac{dda} system.
	\item \textbf{Asteroid System:} Randomly creates asteroid objects which fly across the game field and cause instant death on impact with the player.
	\item \textbf{Collision Detection:} Uses bounding-box collision detection for all game entity interactions.
	\item \textbf{Performance Metrics:} Continuously calculates and exposes player performance metrics (accuracy ratio, kill rate, health ratio, damage rate, survival time) for consumption by the Learning Layer.
\end{itemize}

\subsection{Learning Layer (PPO Model)}
\paragraph{} The Learning Layer implements the \ac{ppo}-based \ac{rl} agent and consists of:
\begin{itemize}
	\item \textbf{Environment Wrapper (\texttt{env\_wrapper\_3d.py}):} A Gymnasium-compatible wrapper that simulates the game environment for training, defining the seven-dimensional observation space and five-action discrete action space.
	\item \textbf{Training Pipeline (\texttt{train\_dda\_3d.py}):} The \ac{ppo} training loop is implemented using PyTorch. The policy network (MLPPolicy3D) uses a shared \ac{mlp} backbone with three hidden layers (256, 256, 128) and separate actor and critic heads.
	\item \textbf{LSTM Extension (\texttt{train\_dda\_3d\_lstm.py}):} Extends the policy network with \ac{lstm} layers to capture temporal patterns in player behavior.
	\item \textbf{Trained Model (\texttt{dda\_ppo\_3d.pth}):} The serialized model weights loaded during live gameplay.
\end{itemize}

\subsection{Decision Layer (Difficulty Adjustment)}
\paragraph{} The Decision Layer bridges the Learning Layer and Gameplay Layer:
\begin{itemize}
	\item At regular intervals during gameplay, the Decision Layer records the current player performance metrics from the Gameplay Layer.
	\item These metrics are formatted as a single observation vector and fed to the trained \ac{ppo} model for inference.
	\item The action taken by the model is mapped onto a set of specific adjustments to game parameters (increase/decrease enemy count and speed).
	\item Gradual adjustments are applied to the Gameplay Layer to ease difficulty changes.
\end{itemize}

\subsection{Analytics Layer (FastAPI + React)}
\paragraph{} The Analytics Layer provides monitoring and visualization capabilities:
\begin{itemize}
	\item \textbf{FastAPI Backend:} Takes the gameplay telemetry data from the game, saves session data, and provides aggregated analytics through \ac{rest}ful \ac{api} endpoints.
	\item \textbf{React Dashboard:} Pulls data from the backend and displays live, up-to-date charts showing player performance metrics, difficulty parameter trends, and session summaries.
\end{itemize}

\begin{figure}[!ht]
	\centering
	\includegraphics[width=0.9\linewidth]{images/system_diagram.png}
	\caption{System architecture diagram showing the four-layer design: Gameplay Layer, Learning Layer, Decision Layer, and Analytics Layer}
	\label{fig:system_architecture}
\end{figure}

\section{Detailed Design}

\subsection{Data Flow Diagram}
\paragraph{} The data flow in the adaptive gameplay system follows a cyclical pattern:

\begin{enumerate}
	\item The player interacts with the game through keyboard inputs.
	\item The Gameplay Layer processes inputs and updates game state, computing player performance metrics.
	\item The Decision Layer periodically collects these metrics and queries the \ac{ppo} model.
	\item The model returns a difficulty adjustment action.
	\item The Decision Layer applies the adjustment to the Gameplay Layer's difficulty parameters.
	\item Simultaneously, gameplay telemetry is sent to the FastAPI backend.
	\item The React dashboard fetches and displays the telemetry data in real time.
\end{enumerate}

\begin{figure}[!ht]
	\centering
	\begin{tikzpicture}[
		node distance=1.5cm and 2cm,
		process/.style={rectangle, minimum width=3cm, minimum height=1cm, text centered, draw=black, fill=blue!10},
		entity/.style={ellipse, minimum width=2.5cm, minimum height=1cm, text centered, draw=black, fill=green!10},
		store/.style={rectangle, minimum width=3cm, minimum height=1cm, text centered, draw=black, fill=yellow!10},
		arrow/.style={thick,->,>=stealth}
		]
		
		\node (player) [entity] {Player};
		\node (gameplay) [process, right=of player] {Gameplay Layer};
		\node (dda) [process, below=of gameplay] {Decision Layer};
		\node (ppo) [process, right=of dda] {Learning Layer};
		\node (api) [store, below=of ppo] {Analytics Backend};
		\node (dashboard) [entity, left=of api] {React Dashboard};
		
		\draw [arrow] (player) -- node[above, font=\footnotesize] {Inputs} (gameplay);
		\draw [arrow] (gameplay) to[bend right] node[above, font=\footnotesize] {Visuals} (player);
		\draw [arrow] (gameplay) -- node[left, font=\footnotesize] {Metrics} (dda);
		\draw [arrow] (dda) -- node[above, font=\footnotesize] {Observation} (ppo);
		\draw [arrow] (ppo) to[bend right] node[below, font=\footnotesize] {Action} (dda);
		\draw [arrow] (dda) to[bend right] node[right, font=\footnotesize] {Adjustments} (gameplay);
		\draw [arrow] (gameplay) to[bend right=45] node[right, near end, font=\footnotesize] {Telemetry} (api);
		\draw [arrow] (api) -- node[above, font=\footnotesize] {JSON} (dashboard);
	\end{tikzpicture}
	\caption{Data Flow Diagram illustrating the cyclical process of telemetry collection, PPO inference, and dynamic game property scaling.}
	\label{fig:dfd}
\end{figure}

\subsection{Class Diagram}
\paragraph{} The system is organized into the following primary classes and modules:

\begin{itemize}
	\item \textbf{Game (game\_3d.py):} The main game class that handles the game loop, rendering, input handling, and entity management. Contains references to Player, Enemy, Asteroid, and Projectile objects.
	\item \textbf{Player:} Manages the player spaceship's position, health, shield, movement speed, and shooting mechanics. Exposes performance metrics for the \ac{dda} system.
	\item \textbf{Enemy:} Represents an enemy ship with configurable speed, health, and attack behavior. Parameters are dynamically adjusted by the \ac{dda} system.
	\item \textbf{Asteroid:} Represents a randomly spawning asteroid object with random spawn position and trajectory.
	\item \textbf{DDAEnv (env\_wrapper\_3d.py):} A Gymnasium-compatible environment class with \texttt{reset()}, \texttt{step()}, and \texttt{render()} methods.
	\item \textbf{MLPPolicy3D (train\_dda\_3d.py):} PyTorch neural network implementing the shared actor-critic policy.
	\item \textbf{PPOTrainer (train\_dda\_3d.py):} Training loop implementing rollout collection, advantage estimation, and policy optimization.
\end{itemize}

\begin{figure}[!ht]
	\centering
	\begin{tikzpicture}[
		class/.style={rectangle, text width=4cm, text centered, draw=black, fill=gray!10, minimum height=1cm},
		arrow/.style={thick,->,>=stealth}
		]
		
		\node (game) [class] {\textbf{Game}\\ \hrule \vspace{2mm} + run()\\ + spawn\_enemies()\\ + update\_difficulty()};
		\node (player) [class, below left=1.5cm and 0.5cm of game] {\textbf{Player}\\ \hrule \vspace{2mm} + health, shield\\ + shots\_hit\\ + update()};
		\node (enemy) [class, below right=1.5cm and 0.5cm of game] {\textbf{Enemy}\\ \hrule \vspace{2mm} + speed, health\\ + attack()\\ + update()};
		\node (env) [class, right=1cm of game] {\textbf{DDAEnv}\\ \hrule \vspace{2mm} + step()\\ + reset()\\ + render()};
		\node (ppo) [class, below=1.5cm of env] {\textbf{MLPPolicy3D}\\ \hrule \vspace{2mm} + actor, critic\\ + forward()};
		
		\draw [arrow] (game) -- node[left, font=\footnotesize] {controls 1} (player);
		\draw [arrow] (game) -- node[right, font=\footnotesize] {manages *} (enemy);
		\draw [arrow] (env) -- node[above, font=\footnotesize] {wraps 1} (game);
		\draw [arrow] (ppo) -- node[right, font=\footnotesize] {interacts 1} (env);
	\end{tikzpicture}
	\caption{Class Diagram depicting the relationships between the Game Engine, Player, DDA Environment, and the PPO MLPPolicy3D component.}
	\label{fig:class_diagram}
\end{figure}

\subsection{Mathematical Modeling of DDA}
\paragraph{} To translate abstract gameplay performance into a standardized format readable by the PPO neural network, the Decision Layer mathematically models the game state. The raw telemetry variables---such as shots fired ($S_f$), shots hit ($S_h$), damage taken ($D_t$), and time elapsed ($T_e$)---exist on vastly different absolute numeric scales. Passing non-normalized integers into a neural network results in severe gradient explosion and policy destabilization.
\paragraph{} Therefore, the Gameplay Layer continuously synthesizes these raw numerical inputs into bounded observation metrics. For example, the accuracy ratio is defined mathematically as $A = S_h / \max(S_f, 1)$, natively binding the metric firmly between $[0, 1]$. Kill rate ($K_r$) and damage rate ($D_r$), which are inherently unbounded and highly dependent on time, are normalized via soft-clipping functions: $K_{norm} = \min(K_r / K_{max}, 1.0)$ where $K_{max}$ represents the theoretical maximum kill throughput per second. The final state observation vector at any discrete timestep $t$ is formulated as the transposed array $O_t = [A, K_{norm}, D_{norm}, H_{ratio}, S_{ratio}, W_{norm}, E_{norm}]^T$. 
\paragraph{} The corresponding action vector $A_t$ generated by the network is a stochastic sample from a discrete probability distribution over five distinctly defined states. Specifically, $A_t \in \{0, 1, 2, 3, 4\}$, representing the relative degree of challenge intensity adjustment. A value of $A_t = 0$ actively maintains the status quo, minimizing gameplay variance. Action subsets $A_t \in \{1, 3\}$ trigger positive deltas (increasing enemy count by $\delta E$ and base velocity by $\delta V$), while $A_t \in \{2, 4\}$ map directly to compensatory negative velocity and count multipliers.

\subsection{Game Engine Subsystems and Physics}
\paragraph{} Constructing the real-time Gameplay Layer required the implementation of several highly specialized subsystems to ensure strict adherence to non-functional latency requirements. The collision detection subsystem utilizes an Axis-Aligned Bounding Box (AABB) intersection algorithm. Unlike computationally complex pixel-perfect polygon masking, AABB operates entirely on hardware-accelerated vector mathematics by exclusively evaluating the rectangular planar coordinates of overlapping sprite boundaries: $(X_1 < X_2 + W_2) \wedge (X_1 + W_1 > X_2) \wedge (Y_1 < Y_2 + H_2) \wedge (Y_1 + H_1 > Y_2)$.
\paragraph{} Furthermore, the engine implements an advanced Object Pooling architecture for projectile and asteroid memory management. Rather than continuously invoking expensive Python object initialization and subsequent garbage collection cycles, dormant visual entities are securely stored in contiguous, pre-allocated memory stacks. When a player actively presses the spacebar, a dormant entity is "woken up," instantly repositioned to the player's coordinate space, and marked active. Upon colliding or leaving the visible viewport, it is deactivated and returned silently to the pool. This protective architectural pattern prevents devastating unhandled memory leaks and ensures absolutely stable $O(1)$ constant time complexity during high-intensity enemy skirmishes.

\subsection{Decision Layer Control Algorithm}
\paragraph{} The continuous interaction between the game engine core and the neural evaluation policy during active play runs entirely asynchronously in order to maintain an unbroken and smooth 60 frames per second standard. Below is the formalized conceptual pseudocode algorithm driving the Decision Layer loops:

\begin{verbatim}
Algorithm 1: Adaptive Decision Loop
Objective: To calculate dynamic game difficulty using PPO RL agent.
Input: Raw telemetric player data metrics
Output: Adjusted Difficulty Parameters (Speed V, Spawn Delay D, Health H)
Steps:
1. Start
2. Initialize Trigger_Interval = 2.0 seconds
3. Loop continuously per game frame:
    4. Update Game Entities and render frame
    5. If elapsed time >= Trigger_Interval:
        6. Read parameters: S_f, S_h, K_r, D_r
        7. Normalize vector O_t to float32 bounds [0, 1]
        8. Compute Logits = Agent_{Forward}(O_t)
        9. Select Action a_t = ArgMax(Logits)
        10. If a_t == INCREASE_HARD:
            11. Assign V = V * 1.10, H = H * 1.05, D = D * 0.90
        12. Else if a_t == DECREASE_EASY:
            13. Assign V = V * 0.90, H = H * 0.95, D = D * 1.10
14. Stop
\end{verbatim}
